% Chapter 4 -- Rubric: PO3 (3.3, 3.4), PO5 (5.1)
\chapter{System Design and Implementation}
\label{ch:design}

\section{Functional Design and Sub-problem Partitioning}
\label{sec:architecture}

\Cref{sec:alternatives} selected a knowledge-graph-backed multi-role agent. This
section sets out how that architecture is partitioned, what each part is
responsible for, and why the boundaries fall where they do.

\subsection{The six sub-problems}

The loop of \Cref{sec:problem-statement} decomposes into six responsibilities.
The partitioning rule throughout is that \emph{a component owns one decision},
so that a change to how that decision is made touches one place.

\begin{enumerate}[label=\textbf{S\arabic*.},leftmargin=3.0em]
  \item \textbf{Knowledge acquisition.} Turn whatever documents exist into
    queryable, individually identified requirements with provenance back to the
    text they came from. Owner: the ingestion pipeline.
  \item \textbf{Test selection.} Decide what is worth testing next, given
    everything currently known. Owner: the planner. This is the exploratory
    problem of \Cref{subsec:primary} and it is where the system's quality is
    bounded.
  \item \textbf{Execution.} Drive a real application toward a stated objective
    and record what happened. Owner: the executor, one per platform.
  \item \textbf{Evidence reduction.} Turn a trajectory into durable,
    deduplicated claims, and decide what the run proved. Owner: the
    investigator.
  \item \textbf{Persistence and retrieval.} Hold the knowledge and answer
    questions about it. Owner: the RAG service, which is the sole writer to the
    database.
  \item \textbf{Operation.} Let a person watch a campaign, understand a verdict
    and take the result away. Owner: the gateway and the dashboard.
\end{enumerate}

\Cref{fig:architecture} shows how these map onto processes. Four run
concurrently: the agent gateway on port 9100, which is the only entry point and
holds the planner; the RAG service on port 9010, which is the only process that
opens a connection to the graph database; the executor, which is a campaign
driver rather than a server; and Neo4j itself on port 7687.

\begin{sidewaysfigure}
  \centering
  \includegraphics[width=0.96\textheight]{01-system-architecture.png}
  \caption[System architecture]{System architecture. Amber is an input, green is
  knowledge, blue is a service, purple is one of the three agent roles, red is a
  language-model call and grey is the device. Every input is optional: the
  system degrades to whatever sources exist rather than failing without them.}
  \label{fig:architecture}
\end{sidewaysfigure}

\subsection{Three decisions that fix the shape}

\paragraph{The graph is the only channel between roles.} The planner, the
executor and the investigator never call one another. Each is a stateless
language-model call that reads from and writes to the same graph, so the
system's memory is entirely in one inspectable place and no conversational state
is carried between rounds. This is what makes ``learning'' a property that can
be examined, rather than an emergent property of a long conversation: the graph
either grew or it did not.

\paragraph{One writer.} Only the RAG service opens a Neo4j connection. The
executor and the gateway reach the graph exclusively through its HTTP interface.
The alternative, letting each process hold a driver, would have distributed the
schema across three codebases and made a constraint change a three-way
migration.

\paragraph{The planner never writes steps.} Its output is an objective in plain
language, not a tap sequence, because it cannot see the live screen. A wrong tap
script is worse than a clear objective, since the executor can recover from an
imprecise goal but not from instructions that do not match the screen in front
of it.

\subsection{The knowledge model}

\Cref{fig:schema} is the schema all six sub-problems share. Its organising idea
is that two graphs coexist in one store: an \emph{a priori} graph of what the
application should do, built from documents, and an \emph{a posteriori} graph of
what it actually does, built from observation. Requirements, chunks and
validation rules belong to the first. UI states, transitions, navigation paths
and execution logs belong to the second. Test cases and findings span them, and
that span is the project's central idea: exploration is steered by the gap
between the two.

\begin{sidewaysfigure}
  \centering
  \includegraphics[width=0.96\textheight]{03-knowledge-graph-schema.png}
  \caption[Knowledge graph schema]{The knowledge graph. Specification-derived
  nodes and observation-derived nodes share one store, joined by the test cases
  and findings that span them.}
  \label{fig:schema}
\end{sidewaysfigure}

\Cref{fig:ingestion} shows S1 in detail, because it is the step that decides
what the rest of the system can ever know. A document is chunked, embedded, and
reduced to individually identified requirements and validation rules, each
linked back to the chunk it was derived from so that any later claim can be
traced to the sentence that justified it. Re-ingesting a changed specification
produces a new version and a drift report rather than overwriting the old one.

\begin{figure}[htbp]
  \centering
  \includegraphics[width=\textwidth]{11-ingestion-document-to-queryable-requirements.png}
  \caption[Ingestion: document to queryable requirements]{Sub-problem S1.
  Prose becomes atomic, identified requirements with provenance back to their
  source chunk, which is what makes requirement coverage reportable at all.}
  \label{fig:ingestion}
\end{figure}

\subsection{Where the standards shaped the partitioning}

Three of the instruments in \Cref{sec:standards} changed the design rather than
merely being satisfied by it.

\emph{Prompt injection (OWASP LLM01).} The application's own on-screen text
reaches a model that also decides the next action. The response is structural:
observed content is labelled as untrusted data in every prompt that carries it,
and the rule appears in the executor, the investigator and the evidence
reviewer rather than in one shared place, because each assembles its own prompt.

\emph{Credential handling (OWASP LLM06).} The test-account secret is passed only
to the executor, which has to type it. The planner receives the account
\emph{role} and nothing else (NFR-D6). This costs a little grounding quality and
removes an entire class of leak, since a secret that never enters the planner's
context cannot appear in a planner trace, a prompt log or the dashboard.

\emph{Accessibility APIs.} Controls are addressed by accessible role, name and
content description rather than by pixel coordinates. That constrains the
executor to what the accessibility layer exposes, and it is what makes the
agent's element-location failures meaningful as an accessibility signal
(\Cref{sec:society}).

\section{Design Refinement and Simulation}
\label{sec:detailed-design}

\subsection{One round, end to end}

\Cref{fig:round} is the control flow of a single test, and it is the view to
read first: it shows every component, every artefact and the point at which the
graph grows.

\begin{figure}[htbp]
  \centering
  \includegraphics[width=0.96\textwidth]{02-one-test-round-end-to-end.png}
  \caption[One test round, end to end]{One test round. The planner reads the
  graph, proposes, and is validated before anything reaches a device. The
  executor drives up to 50 steps. The investigator reduces the trajectory to
  findings, and is given the findings already known for those screens so that a
  repeat observation reinforces an existing finding instead of creating a
  duplicate. The next round starts from the larger graph.}
  \label{fig:round}
\end{figure}

\subsection{The refinement that mattered most: retrieval to investigation}

The planner was built twice, and the second design exists because the first had
a structural defect rather than a tuning problem.

In the original \emph{pipeline} planner each knowledge source returned a block
carrying both a body of text and a one-line note. The text went into a bucket
that was not read again until generation; only the note reached the next
retrieval decision. The planner therefore chose what to fetch next
\emph{without having read what it had just fetched}, and everything was
concatenated into a single prompt of roughly 19{,}000 characters at the end.
\Cref{fig:tools-vs-sources} contrasts the two.
\Cref{fig:pipeline} is the pipeline planner itself: a five-node state machine of
which only three nodes call a model.

\begin{figure}[htbp]
  \centering
  \includegraphics[height=0.70\textheight,keepaspectratio]{05-planner-pipeline-mode-default.png}
  \caption[The pipeline planner]{The pipeline planner, still the default. The
  defect is not in any single node but in the edge between \texttt{planner\_step}
  and \texttt{execute\_retrieval}: only a one-line note returns, so the next
  retrieval decision is taken without having read the previous result.}
  \label{fig:pipeline}
\end{figure}


\begin{figure}[htbp]
  \centering
  \includegraphics[height=0.72\textheight,keepaspectratio]{13-why-tools-replaced-retrieval-sources.png}
  \caption[Retrieval sources versus tools]{Why tools replaced retrieval sources.
  Above, the model sees no content until the single generation call. Below, each
  result enters the conversation, so the second question can depend on the answer
  to the first.}
  \label{fig:tools-vs-sources}
\end{figure}

Three symptoms followed from that one cause, and each had been patched
separately before the cause was identified: screen names that did not exist,
duplicate detection that could only run after generation, and requirement
identifiers the model had invented. All three are the same missing capability,
namely that the planner could not check a claim against the graph while forming
it.

The refinement was to give the planner tools instead of sources.
\Cref{fig:nine-tools} groups the nine tools by the question each answers, and
the grouping is the architecture in miniature: what \emph{should} exist, what
the application \emph{actually} has, what is \emph{already known}, and therefore
where the gap is.

\begin{figure}[htbp]
  \centering
  \includegraphics[height=0.74\textheight,keepaspectratio]{14-the-nine-tools-as-an-investigation.png}
  \caption[The nine planner tools]{The nine tools as an investigation. A tool
  whose knowledge source is disabled or empty is never registered, so an absent
  source cannot mislead the model with a blank answer.}
  \label{fig:nine-tools}
\end{figure}

\paragraph{Measured effect.} Under the pipeline planner, the screen a proposal
named matched a screen the executor had actually observed about 16\% of the
time, and the goal text had to hedge by calling the name a lead rather than a
fact while the executor spent its budget searching for screens that did not
exist. Under the tool-calling planner, across eight test cases there was one
rejection in total, corrected on the following turn. Both planners remain
selectable by one setting, \texttt{PLANNER\_MODE}, and produce the same output
contract, so a campaign can be run on each and compared without touching the
executor. The default remains \texttt{pipeline} until that comparison has been
made over a full campaign.

\subsection{Design calculations and budgets}

Every bound in the system is explicit, because an exploratory agent has no
natural stopping point (difficulty D2) and will otherwise consume whatever it is
given. \Cref{tab:budgets} lists them.

\begingroup
\small
\setlength{\tabcolsep}{5pt}
\begin{xltabular}{\textwidth}{@{}>{\hsize=0.80\hsize}Y>{\hsize=0.44\hsize}Y>{\hsize=1.76\hsize}Y@{}}
  \caption{Explicit bounds in the design, and the reasoning behind each}
  \label{tab:budgets}\\
  \toprule
  Bound & Value & Why it exists \\
  \midrule
  \endfirsthead
  \caption[]{Explicit bounds in the design \emph{(continued)}}\\
  \toprule
  Bound & Value & Why it exists \\
  \midrule
  \endhead
  \midrule
  \multicolumn{3}{r@{}}{\footnotesize\itshape continued on the next page}\\
  \endfoot
  \bottomrule
  \endlastfoot

  Device actions per test & 50 & The executor's interaction budget. Exhausting it yields no evidence about the application, so the planner is told the number and is warned when recent tests repeatedly hit it. \\
  \addlinespace
  Wall clock per test & 900\,s & Bounds a single execution including model latency and recovery. \\
  \addlinespace
  Planner turns & 12 & Ceiling on the tool loop. \\
  \addlinespace
  Forced proposal after & 6 turns & From this turn the model is told to propose and the tool choice is pinned, because a model left free will investigate to the ceiling and never commit. Observed on the first live run: ten turns of tool calls and no proposal. \\
  \addlinespace
  Proposal rejections & 3 & After three, the best effort is accepted and a degradation is recorded, so a planning round never returns nothing. \\
  \addlinespace
  Tool result caps & 3000 / 2500 / 1200 / 2000 characters & Requirements, findings, one screen, and any other result. There is no global prompt budget in the tool planner; each tool is responsible for not flooding the context on its own. \\
  \addlinespace
  Finding attempts & 3 & After three probes an open question auto-closes as inconclusive, which is a result for a human rather than an endless retry. \\
  \addlinespace
  Knowledge half-life & 90 days & Decays the weight of stale observations in retrieval so the agent adapts when the application changes. \\
  \bottomrule
\end{xltabular}
\endgroup

\subsection{State identity: the analytical core}

The single most consequential calculation in the system is not a budget but a
similarity decision. Every observation must be classified as a screen already in
the application model or as a new one. Merging too eagerly collapses the map;
splitting too eagerly destroys it. One wizard, whose form fields carried their
current values in the same attribute that carries a button's label, produced
nine distinct states before the cascade in \Cref{fig:state-identity} was
designed.

\begin{figure}[htbp]
  \centering
  \includegraphics[height=0.82\textheight,keepaspectratio]{17-state-identity-is-this-a-screen-we-have-seen.png}
  \caption[State identity cascade]{The five-stage state-identity cascade,
  cheapest test first. Each stage was added in response to a specific way the
  previous one failed, and every threshold is a configurable setting.}
  \label{fig:state-identity}
\end{figure}

The stages are ordered by cost. An exact signature match is a string comparison.
Structural Jaccard tolerates dynamic chrome such as counters and timestamps.
Containment recognises a screen captured before its late-painting chrome, which
is a strict subset of the same screen fully rendered and which Jaccard therefore
scores as different; it is guarded by a minimum control count and a maximum size
ratio, because a one-control blank screen is contained in everything. Skeleton
containment compares only resource identifier and class, asking whether the same
widgets are present regardless of their values, with a floor on full containment
so that two genuinely different screens built from similar widget types stay
apart. The perceptual hash is the last resort and the only signal available when
the accessibility tree is too thin to describe the screen at all.

\subsection{Attribution: keeping three kinds of failure apart}

A run that does not pass is ambiguous, and conflating the three explanations
inflates the reported defect count, which is the most expensive error this
system can make. \Cref{fig:attribution} shows the classification.

\begin{figure}[htbp]
  \centering
  \includegraphics[width=\textwidth]{09-failure-attribution.png}
  \caption[Failure attribution]{Failure attribution. The three sets have exactly
  one definition in the source tree. Only the first is evidence about the
  application.}
  \label{fig:attribution}
\end{figure}

The design decision worth noting is that the taxonomy is defined once, as three
frozen sets in the settings module, and imported everywhere it is needed. An
earlier version defined it in two places, and the two drifted: autonomy was
reported as 100\% when it was 67\%.

\subsection{Simulation and verification of the design}

Three mechanisms verify the design without consuming a device or a model budget.
A \emph{simulator runner} exercises the full planning and logging loop with a
stubbed executor, so the graph writes, the coverage calculation and the
attribution path can be tested without a handset. A set of \emph{configuration
guards} asserts the invariants that the design depends on, including that the
attribution sets are disjoint and that no application-specific vocabulary has
entered the code. And the \emph{state-identity} and \emph{observation} modules
are pure functions over recorded inputs, so the cascade above is tested on fixed
key sets rather than on a live device.

\section{Application of Tools}
\label{sec:tools}

\Cref{fig:topology} is the deployment view of the same system: four processes on
one host, one device over the Android Debug Bridge, and exactly one external
dependency.

\begin{figure}[htbp]
  \centering
  \includegraphics[width=\textwidth]{12-runtime-topology.png}
  \caption[Runtime topology]{Runtime topology. The model provider is the only
  component outside the developer's machine.}
  \label{fig:topology}
\end{figure}

\Cref{tab:tools} states why each significant tool was chosen and what it costs.

\begingroup
\small
\setlength{\tabcolsep}{5pt}
\begin{xltabular}{\textwidth}{@{}>{\hsize=0.54\hsize}Y>{\hsize=1.18\hsize}Y>{\hsize=1.28\hsize}Y@{}}
  \caption{Tool selection and the limitations accepted with each}
  \label{tab:tools}\\
  \toprule
  Tool & Why it was selected & Limitation accepted \\
  \midrule
  \endfirsthead
  \caption[]{Tool selection and limitations \emph{(continued)}}\\
  \toprule
  Tool & Why it was selected & Limitation accepted \\
  \midrule
  \endhead
  \midrule
  \multicolumn{3}{r@{}}{\footnotesize\itshape continued on the next page}\\
  \endfoot
  \bottomrule
  \endlastfoot

  Neo4j & The payload is the relationships: a test covers a requirement, a log belongs to a test, a navigation node leads to a UI state, a defect concentrates in an area. These are traversals, and a property graph answers them directly while also hosting the vector indexes dense retrieval needs. &
  A second service to run and back up. Cypher is another language for the team to learn, and the community edition is GPL-licensed, which constrains redistribution. \\
  \addlinespace
  mobilerun 0.6.8 & Drives Android through the accessibility service with a companion on-device application that also supplies an input method, which is what makes text entry into multi-field forms reliable. Exposes an event stream the executor subscribes to for observed states. &
  A young dependency whose API shape has changed between versions. Its accessibility client conflicts with \texttt{uiautomator} if both are registered. \\
  \addlinespace
  Playwright & The web counterpart: an accessibility-annotated snapshot rather than raw DOM, which lets the same planner address controls by role and name on both platforms. &
  Browser automation cannot see native dialogs outside the page, and file-chooser handling is an explicit special case. \\
  \addlinespace
  LangGraph & Expresses the pipeline planner as an explicit state machine rather than a chain of calls, so each node is separately testable and the flow is inspectable. &
  Adds a framework dependency for what is, in the tool planner, a loop that could be written directly. \\
  \addlinespace
  OpenRouter & One interface to several model providers, which made the tier downgrades of \Cref{sec:economics} a configuration change rather than an integration. &
  The same model is served by different providers with order-of-magnitude differences in latency and token usage, which adds noise to any cost or timing measurement. \\
  \addlinespace
  fastembed with \modpath{BAAI/bge-small-en-v1.5} & Runs embeddings locally through ONNX with no heavyweight tensor framework and no per-call charge, at 384 dimensions. &
  An English-centric model used partly on Bengali interface text, which is a known weakness of the retrieval path rather than a solved problem. \\
  \addlinespace
  FastAPI and Uvicorn & Typed request and response schemas that double as the interface documentation, with the validation the graph writes depend on. &
  None material at this scale. \\
  \addlinespace
  React with Vite & The operator surface builds to a single self-contained file the gateway serves directly, with no separate deployment. &
  A JavaScript toolchain in an otherwise Python project. \\
  \bottomrule
\end{xltabular}
\endgroup

\section{Implementation}
\label{sec:implementation}

\subsection{The executor in detail}

\Cref{fig:executor} is the most operationally dense part of the system: how a
goal becomes device actions, what is observed, and every artefact the run leaves
behind.

\begin{sidewaysfigure}
  \centering
  \includegraphics[width=0.96\textheight]{16-the-executor-in-detail.png}
  \caption[The executor in detail]{The executor. Note the two feedback paths to
  the operator: the device agent's log is teed to disk and streamed live, and
  the per-step trajectory is replayed on demand after the run.}
  \label{fig:executor}
\end{sidewaysfigure}

Two implementation decisions are worth stating because neither is obvious from
the architecture.

\paragraph{Observation is retried, and its failure is recorded.} The
accessibility tree is not always available at the instant it is requested,
particularly while a screen is still painting. The observation helper retries
three times before giving up, and a give-up is written to the degradation
register rather than being silently tolerated, because an observation that
quietly returned nothing would produce a state map with holes in it and no
indication that holes existed.

\paragraph{There is no in-loop livelock detector on Android.} An earlier version
had one. It was removed because cancelling an agent mid-run left that agent
executing alongside the next test, which corrupted both. The capability was
therefore moved rather than restored: over-scoping now surfaces as a step-limit
failure, and the \emph{planner} corrects it from the recent-runs block, which
tells it when consecutive tests have exhausted their budget. The web executor,
whose cancellation is a browser close and therefore safe, does keep two in-loop
guards. This asymmetry is deliberate and is the clearest example in the project
of a capability being relocated to where it can be implemented safely.

\subsection{The investigator}

\Cref{fig:investigator} is sub-problem S4: how a trajectory of device steps
becomes knowledge that survives the run. The design decision that matters is
that the investigator is given the findings already recorded for the screens the
run touched. A repeated observation therefore increments an existing finding's
independent-observation count instead of creating a near-duplicate, which is
what allows a claim seen three times across three separate evaluations to be
distinguished from a claim asserted once.

\begin{figure}[htbp]
  \centering
  \includegraphics[width=\textwidth]{06-investigator-trajectory-to-knowledge.png}
  \caption[The investigator]{Sub-problem S4. A trajectory becomes typed,
  deduplicated findings, and the verdict is separated from the findings so that
  a failed run can still contribute knowledge.}
  \label{fig:investigator}
\end{figure}

\subsection{The operator surface}

\Cref{fig:dash-logs} shows the two live log panes the executor feeds: the device
agent's own reasoning on the left, streamed from the teed log file, and the
planner's node-by-node trace on the right. They stream independently, so a stall
is immediately attributable to one side or the other.

\begin{figure}[htbp]
  \centering
  \includegraphics[width=\textwidth]{dashboard-live-logs.jpg}
  \caption[Dashboard: live logs]{The dashboard during a campaign. Device and
  planner logs stream independently, with the generated test cases and the
  findings so far below.}
  \label{fig:dash-logs}
\end{figure}

\Cref{fig:dash-model} shows the application model the agent has built for
itself, as a navigable graph of screens and transitions with the screenshot of
any selected state.

\begin{figure}[htbp]
  \centering
  \includegraphics[width=\textwidth]{dashboard-app-model.jpg}
  \caption[Dashboard: the live application model]{The Live App Model panel. Each
  node is a UI state the agent reached, deduplicated by the cascade of
  \Cref{fig:state-identity}; each edge is a transition it took.}
  \label{fig:dash-model}
\end{figure}

\guidance{TODO (team): these two captures are from an earlier campaign against a
contacts application and are included here as evidence that the operator surface
exists, which is what this section claims. Do not reuse them in
\Cref{ch:evaluation}: their totals are from a different application and would
contradict the campaign reported there. Take fresh captures during a current run
if \Cref{ch:evaluation} needs a live-run figure.}

\subsection{Code organisation}

\Cref{tab:modules} maps each sub-problem of \Cref{sec:architecture} to the code
that implements it.

\begingroup
\small
\setlength{\tabcolsep}{5pt}
\begin{xltabular}{\textwidth}{@{}>{\hsize=0.50\hsize}Y>{\hsize=0.96\hsize}Y>{\hsize=1.54\hsize}Y@{}}
  \caption{Sub-problems and the modules that implement them}
  \label{tab:modules}\\
  \toprule
  Sub-problem & Module & Responsibility \\
  \midrule
  \endfirsthead
  \caption[]{Sub-problems and modules \emph{(continued)}}\\
  \toprule
  Sub-problem & Module & Responsibility \\
  \midrule
  \endhead
  \midrule
  \multicolumn{3}{r@{}}{\footnotesize\itshape continued on the next page}\\
  \endfoot
  \bottomrule
  \endlastfoot

  S1 Knowledge acquisition & \modpath{ingestion/} & Document loading, multi-sample extraction with a judge model, and the canonical UI representation shared by both platforms. \\
  \addlinespace
  S2 Test selection & \modpath{planner/} & Both planner modes (\modpath{agent_loop.py} and \modpath{langgraph_agent.py}), the nine tools, the validating proposal gate, coverage and the exploration directive, and the pluggable source registry. \\
  \addlinespace
  S3 Execution & \modpath{clients/executor_runner.py}, \modpath{web_player/} & Android and web execution, each with its own observation, action and failure-classification code, sharing one attribution vocabulary. \\
  \addlinespace
  S4 Evidence reduction & \modpath{rag_api/findings.py} & Trajectory to typed, deduplicated findings with an independent-observation count and a lifecycle. \\
  \addlinespace
  S5 Persistence and retrieval & \modpath{rag_api/} & The sole database owner: schema, constraints, vector indexes, retrieval, the state-identity cascade, learning, risk, metrics and anomalies. \\
  \addlinespace
  S6 Operation & \modpath{gateway/}, \modpath{dashboard-react/}, \modpath{targets/} & The single entry point, the operator dashboard, the exportable run report, and the validated target profiles that make a new application a configuration change. \\
  \bottomrule
\end{xltabular}
\endgroup

\subsection{Development process}

Work was distributed by subsystem, each member holding a branch and integrating
through the shared default branch. The repository carries an automated test
suite of six modules run by a single entry point; at the time of writing 124
checks pass without any external service, with the findings module additionally
requiring a live database. The suite is deliberately weighted toward the parts
of the system that are pure functions over recorded inputs, namely state
identity, observation normalisation and the configuration invariants, because
those are the parts where a regression would be silent.

Two practices were adopted after they proved necessary rather than in advance.
Every degradation is recorded and readable across process boundaries, so a
component that falls back cannot do so quietly. And application-agnosticism is
checked by an automated sweep of the source tree rather than by review, after a
manual audit declared the system generic while application-specific vocabulary
was still embedded in the executor and in prompt templates.

\guidance{TODO (team): add the repository URL, the branch and review convention
actually used, and whether any continuous-integration workflow runs the test
suite on push. Keep this consistent with \Cref{sec:participation}.}

\section{Deployment}
\label{sec:deployment}

\subsection{Target environment}

The system is deployed as four processes on a single developer machine, as
\Cref{fig:topology} shows. A launch script starts the RAG service and the
gateway under Uvicorn, waits for both health endpoints, and optionally starts the
executor campaign loop. Neo4j runs locally. The Android device is a physical
handset or an emulator attached over the Android Debug Bridge, carrying the
companion application that provides the accessibility service and the input
method. Model inference is the only component reached over the network.

\subsection{Configuration and release}

Configuration is layered deliberately. Process-wide defaults live in a settings
module read once at import, which is why a campaign is launched as a subprocess
rather than in-process: the settings object is global, so running a second
target inside a live server would leak that target's configuration into every
other request. Per-application configuration is a validated JSON target profile,
which the dashboard can list, validate, save, ingest and run, so adding an
application is a data change rather than a code change. Secrets live only in an
environment file that is excluded from version control.

\subsection{Operation and monitoring}

A campaign is operated from the dashboard: the live session, the growing
application model, both log streams and the planner trace are visible while it
runs, and a run report exports as a PDF afterwards. Three artefact streams are
written to disk for after the fact: structured application logs, one trajectory
per execution, and the degradation register. The graph itself is the durable
record and is the thing that must be backed up; snapshots are held outside
version control, under the same controls as the test credentials, for the
reasons given in \Cref{sec:sustainability}.

\subsection{Scope of the deployment claim}

This is a local, developer-operated deployment on a trusted host, which is
constraint C7 and is stated as a boundary rather than discovered as a defect.
The services bind to a host-local interface by default and support key-based
authentication; deploying beyond a single trusted host would require uniform
authentication across every diagnostic handler and containment checks on
ingestion paths, which is identified in \Cref{sec:standards} as a prerequisite
and was not in scope. The computing resources this deployment consumes, and
their environmental cost, are assessed in \Cref{sec:sustainability}.
